% Part: first-order-logic
% Chapter: completeness
% Section: complete-sets

% Definition of complete consistent sets. Properties of complete sets required
% for completeness proved are in provability.tex in the
% chapter on the proof system used.

\documentclass[../../../include/open-logic-section]{subfiles}

\begin{document}

\iftag{FOL}
      {\olfileid[hi]{fol}{com}{ccs}}
      {\olfileid[hi]{pl}{com}{ccs}}
      
\olsection{\usetoken{P}{sentence} के पूर्ण संगत समुच्चय}

\begin{defn}[पूर्ण समुच्चय]
\ollabel{def:complete-set} !!{sentence}s का समुच्चय~$\Gamma$
\emph{!!{complete}} तभी और केवल तभी है जब किसी भी !!{sentence}~$!A$ के लिए
या तो $!A \in
\Gamma$ हो या $\lnot !A \in \Gamma$ हो।
\end{defn}

\begin{explain}
वाक्यों के !!^{complete} समुच्चय किसी प्रश्न को अनुत्तरित नहीं छोड़ते। किसी
भी !!{sentence}~$!A$ के लिए $\Gamma$ यह ``कहता'' है कि $!A$ सत्य है या
असत्य। !!{complete} समुच्चयों का महत्त्व पूर्णता प्रमेय के प्रमाण से आगे तक
जाता है। उदाहरण के लिए, कोई सिद्धांत जो !!{complete} और अभिगृहीतीकरणीय हो,
सदा निर्णेय होता है।
\end{explain}

\begin{explain}
पूर्णता के प्रमाण में !!^{complete} संगत समुच्चय महत्त्वपूर्ण हैं, क्योंकि
हम गारंटी दे सकते हैं कि !!{sentence}s का प्रत्येक संगत समुच्चय~$\Gamma$
किसी !!a{complete} संगत समुच्चय~$\Gamma^*$ में समाहित है। प्रत्येक
!!^a{complete} संगत समुच्चय में, हर !!{sentence}~$!A$ के लिए, या तो $!A$ या
उसका निषेध $\lnot !A$ होता है, पर दोनों नहीं। विशेषतः यह
\iftag{FOL}{परमाण्विक !!{sentence}s}{प्रतिज्ञप्ति चरों} के लिए सत्य है; अतः
!!a{complete} संगत समुच्चय\iftag{FOL}{ से, ऐसी भाषा में जिसे
!!{constant}s द्वारा उपयुक्त रूप से विस्तारित किया गया हो}{}, हम ऐसी
\iftag{FOL}{!!a{structure}}{!!a{valuation}} बना सकते हैं जिसमें
\iftag{FOL}{!!{predicate}s की व्याख्या}{प्रतिज्ञप्ति चरों को दिया गया
सत्य-मूल्य} इस आधार पर परिभाषित होता है कि कौन-से
\iftag{FOL}{परमाण्विक !!{sentence}s}{!!{propositional variable}s}
समुच्चय~$\Gamma^*$ में हैं। फिर दिखाया जा सकता है कि यह
\iftag{FOL}{!!{structure}}{!!{valuation}},~$\Gamma^*$ के सभी !!{sentence}s
को (और इसलिए~$\Gamma$ के सभी वाक्यों को भी) सत्य बनाती है। इस अंतिम तथ्य के
प्रमाण के लिए यह आवश्यक है कि $\lnot !A \in \Gamma^*$ तभी और केवल तभी जब
$!A \notin \Gamma^*$ हो, $(!A \lor !B) \in
\Gamma^*$ तभी और केवल तभी जब $!A \in \Gamma^*$ या $!B \in \Gamma^*$ हो,
इत्यादि।
\end{explain}

आगे हम प्रायः $\Proves$ की स्वतुल्यता, एकदिष्टता और संक्रामकता के गुणों का
निःशब्द उपयोग करेंगे (देखें
\iftag{FOL}{%
  \tagrefs{prfSC/{fol:seq:ptn:sec},prfND/{fol:ntd:ptn:sec},prfAX/{fol:axd:ptn:sec},prfTab/{fol:tab:ptn:sec}}}{%
  \tagrefs{prfSC/{pl:seq:ptn:sec},prfND/{pl:ntd:ptn:sec},prfAX/{pl:axd:ptn:sec},prfTab/{pl:tab:ptn:sec}}})।

\begin{prop}
\ollabel{prop:ccs}
मान लें कि $\Gamma$ !!{complete} और संगत है। तब:
\begin{enumerate}
\item \ollabel{prop:ccs-prov-in} यदि $\Gamma \Proves !A$, तो $!A \in
  \Gamma$।

\tagitem{prvAnd}{\ollabel{prop:ccs-and} $!A \land !B \in \Gamma$
  तभी और केवल तभी जब $!A \in \Gamma$ और $!B \in \Gamma$ दोनों हों।}{}

\tagitem{prvOr}{\ollabel{prop:ccs-or} $!A \lor !B \in \Gamma$ तभी और केवल
  तभी जब या तो $!A \in \Gamma$ या $!B \in \Gamma$ हो।}{}

\tagitem{prvIf}{\ollabel{prop:ccs-if} $!A \lif !B \in \Gamma$ तभी और केवल
  तभी जब या तो $!A \notin \Gamma$ या $!B \in \Gamma$ हो।}{}
\end{enumerate}
\end{prop}

\begin{proof}
निम्न सभी बातों के लिए मान लें कि $\Gamma$ !!{complete} और संगत है।
\begin{enumerate}
\item यदि $\Gamma \Proves !A$, तो $!A \in \Gamma$।

मान लें कि $\Gamma \Proves !A$। इसके विपरीत मान लें कि $!A
\notin \Gamma$। चूँकि $\Gamma$ !!{complete} है, $\lnot !A \in \Gamma$।
\iftag{FOL}{%
  \tagrefs{prfSC/{fol:seq:prv:prop:explicit-inc},
    prfND/{fol:ntd:prv:prop:explicit-inc},
    prfAX/{fol:axd:prv:prop:explicit-inc},
    prfTab/{fol:tab:prv:prop:explicit-inc}}}{%
  \tagrefs{prfSC/{pl:seq:prv:prop:explicit-inc},
    prfND/{pl:ntd:prv:prop:explicit-inc},
    prfAX/{pl:axd:prv:prop:explicit-inc},
    prfTab/{pl:tab:prv:prop:explicit-inc}}} के अनुसार
$\Gamma$ असंगत है। यह इस धारणा का खंडन करता है कि $\Gamma$ संगत है। अतः
$!A \notin
\Gamma$ नहीं हो सकता, इसलिए $!A \in \Gamma$।

\tagitem{defAnd}{}{%
\iftag{probAnd}{अभ्यास।}{%
$!A \land !B \in \Gamma$ तभी और केवल तभी जब $!A \in \Gamma$ और $!B \in \Gamma$
दोनों हों:

अग्र दिशा के लिए मान लें कि $!A \land !B \in \Gamma$। तब
\iftag{FOL}{%
  \tagrefs{prfSC/{fol:seq:ppr:prop:provability-land},
    prfND/{fol:ntd:ppr:prop:provability-land},
    prfAX/{fol:axd:ppr:prop:provability-land},
    prfTab/{fol:tab:ppr:prop:provability-land}}}{%
  \tagrefs{prfSC/{pl:seq:ppr:prop:provability-land},
    prfND/{pl:ntd:ppr:prop:provability-land},
    prfAX/{pl:axd:ppr:prop:provability-land},
    prfTab/{pl:tab:ppr:prop:provability-land}}} के पद~(1) से
$\Gamma \Proves !A$ और $\Gamma \Proves !B$। \olref{prop:ccs-prov-in} से
$!A \in \Gamma$ और $!B \in \Gamma$ मिलते हैं, जैसा अपेक्षित था।

विपरीत दिशा के लिए $!A \in \Gamma$ और $!B \in
\Gamma$ मान लें।
\iftag{FOL}{%
  \tagrefs{prfSC/{fol:seq:ppr:prop:provability-land},
    prfND/{fol:ntd:ppr:prop:provability-land},
    prfAX/{fol:axd:ppr:prop:provability-land},
    prfTab/{fol:tab:ppr:prop:provability-land}}}{%
  \tagrefs{prfSC/{pl:seq:ppr:prop:provability-land},
    prfND/{pl:ntd:ppr:prop:provability-land},
    prfAX/{pl:axd:ppr:prop:provability-land},
    prfTab/{pl:tab:ppr:prop:provability-land}}} के पद~(2) से
$\Gamma \Proves !A \land !B$। \olref{prop:ccs-prov-in} से $!A \land
!B \in \Gamma$।}}

\tagitem{defOr}{}{%
\iftag{probOr}{अभ्यास।}{%
पहले हम दिखाते हैं कि यदि $!A \lor !B \in \Gamma$, तो या तो $!A \in
\Gamma$ या $!B \in \Gamma$। मान लें कि $!A \lor !B \in \Gamma$ पर $!A
\notin \Gamma$ और $!B \notin \Gamma$। चूँकि $\Gamma$
!!{complete} है, $\lnot !A \in \Gamma$ और $\lnot !B \in \Gamma$।
\iftag{FOL}{%
  \tagrefs{prfSC/{fol:seq:ppr:prop:provability-lor},
    prfND/{fol:ntd:ppr:prop:provability-lor},
    prfAX/{fol:axd:ppr:prop:provability-lor},
    prfTab/{fol:tab:ppr:prop:provability-lor}}}{%
  \tagrefs{prfSC/{pl:seq:ppr:prop:provability-lor},
    prfND/{pl:ntd:ppr:prop:provability-lor},
    prfAX/{pl:axd:ppr:prop:provability-lor},
    prfTab/{pl:tab:ppr:prop:provability-lor}}} के पद (1) से $\Gamma$ असंगत
है, जो विरोधाभास है। अतः या तो $!A
\in \Gamma$ या $!B \in \Gamma$।

विपरीत दिशा के लिए मान लें कि $!A \in \Gamma$ या $!B \in
\Gamma$।
\iftag{FOL}{%
  \tagrefs{prfSC/{fol:seq:ppr:prop:provability-lor},
    prfND/{fol:ntd:ppr:prop:provability-lor},
    prfAX/{fol:axd:ppr:prop:provability-lor},
    prfTab/{fol:tab:ppr:prop:provability-lor}}}{%
  \tagrefs{prfSC/{pl:seq:ppr:prop:provability-lor},
    prfND/{pl:ntd:ppr:prop:provability-lor},
    prfAX/{pl:axd:ppr:prop:provability-lor},
    prfTab/{pl:tab:ppr:prop:provability-lor}}} के पद (2) से
$\Gamma \Proves !A \lor !B$। \olref{prop:ccs-prov-in} से $!A \lor
!B \in \Gamma$, जैसा अपेक्षित था।}}

\tagitem{defIf}{}{%
\iftag{probIf}{अभ्यास।}{%
अग्र दिशा के लिए मान लें कि $!A \lif !B \in \Gamma$, और इसके विपरीत मान
लें कि $!A \in \Gamma$ तथा $!B \notin \Gamma$। इन धारणाओं पर
$!A \lif !B \in \Gamma$ और $!A \in \Gamma$।
\iftag{FOL}{%
  \tagrefs{prfSC/{fol:seq:ppr:prop:provability-lif},
    prfND/{fol:ntd:ppr:prop:provability-lif},
    prfAX/{fol:axd:ppr:prop:provability-lif},
    prfTab/{fol:tab:ppr:prop:provability-lif}}}{%
  \tagrefs{prfSC/{pl:seq:ppr:prop:provability-lif},
    prfND/{pl:ntd:ppr:prop:provability-lif},
    prfAX/{pl:axd:ppr:prop:provability-lif},
    prfTab/{pl:tab:ppr:prop:provability-lif}}} के पद (1) से
$\Gamma \Proves !B$। किंतु तब \olref{prop:ccs-prov-in} से $!B \in
\Gamma$, जो $!B \notin \Gamma$ वाली धारणा का खंडन करता है।

विपरीत दिशा के लिए पहले उस स्थिति पर विचार करें जिसमें $!A \notin
\Gamma$। चूँकि $\Gamma$ !!{complete} है, $\lnot !A \in \Gamma$।
\iftag{FOL}{%
  \tagrefs{prfSC/{fol:seq:ppr:prop:provability-lif},
    prfND/{fol:ntd:ppr:prop:provability-lif},
    prfAX/{fol:axd:ppr:prop:provability-lif},
    prfTab/{fol:tab:ppr:prop:provability-lif}}}{%
  \tagrefs{prfSC/{pl:seq:ppr:prop:provability-lif},
    prfND/{pl:ntd:ppr:prop:provability-lif},
    prfAX/{pl:axd:ppr:prop:provability-lif},
    prfTab/{pl:tab:ppr:prop:provability-lif}}} के पद (2) से
$\Gamma \Proves !A \lif !B$। पुनः \olref{prop:ccs-prov-in} से
$!A \lif !B \in \Gamma$ मिलता है, जैसा अपेक्षित था।

अब उस स्थिति पर विचार करें जिसमें $!B \in \Gamma$।
\iftag{FOL}{%
  \tagrefs{prfSC/{fol:seq:ppr:prop:provability-lif},
    prfND/{fol:ntd:ppr:prop:provability-lif},
    prfTab/{fol:tab:ppr:prop:provability-lif},
    prfAX/{fol:axd:ppr:prop:provability-lif}}}{%
  \tagrefs{prfSC/{pl:seq:ppr:prop:provability-lif},
    prfND/{pl:ntd:ppr:prop:provability-lif},
    prfTab/{pl:tab:ppr:prop:provability-lif},
    prfAX/{pl:axd:ppr:prop:provability-lif}}} के पद (2) से फिर
$\Gamma \Proves !A \lif !B$। \olref{prop:ccs-prov-in} से
$!A \lif !B \in \Gamma$।}}
\end{enumerate}
\end{proof}

\tagprob{FOL}
\begin{prob}
\olref[fol][com][ccs]{prop:ccs} का प्रमाण पूरा कीजिए।
\end{prob}
\tagendprob

\tagprob{notFOL}
\begin{prob}
\olref[pl][com][ccs]{prop:ccs} का प्रमाण पूरा कीजिए।
\end{prob}
\tagendprob

\end{document}
